\subsection{Proof of Lemma~\ref{lem:round-j}: Rounding Subgraphs with Large Shift Potential}
\label{sec:low-ent-j}
In this section, we will show that when the shift-partition potential $\Phi^G_{\beta,\nu}$ (Definition~\ref{def:shift-partition}) is large with respect to a pseudodistribution $\cD$ (with certain nice properties), then the Condition\&Round Algorithm  (Algorithm~\ref{alg:low-ent-restated}) succeeds in returning a good assignment for the unique games instance. The proofs in this section follow along the lines of the analysis of Condition\&Round given in~\cite{BBKSS}, albeit instead of independence between $X,X'$ we only have approximate local independence. Since we will always be working with the graph $G$ we will henceforth drop the superscript $G$ from the shift-partition potential $\Phi$.
We will prove the following theorem in this section:

\begin{lemma}[Restatement of Lemma~\ref{lem:round-j}]\label{thm:round-j}
Let $I = (G, \Pi)$ be an affine instance of Unique Games over the alphabet $\Sigma$ and $\mu_1 \times \mu_2$ be a degree $D+\widetilde{O}(1/\nu)$ pseudodistribution over assignments $(X,X')$ to $I$. Let $E(X,X')$ be a polynomial such that $\cA_I(X,X') \vdash_D E(X,X') \in [0,1]$. Suppose we have that: 
\[\Pr_{u,v \sim V(G): u \neq v}[TV((Y_{u,v},Y'_{u,v})|E , (Y_{u,v},Y'_{u,v})) > \delta] \leq \zeta,\]
where the distribution $(Y_{u,v},Y'_{u,v})|E$ refers to the joint distribution on these variables defined by the collection of local distributions $\cL(\mu_1 \times \mu_2|E, p_{\beta,\nu})$ and similarly $(Y_{u,v},Y'_{u,v})$ refers to the distribution defined by $\cL(\mu_1 \times \mu_2,p_{\beta,\nu})$ (Definition~\ref{defn:local-dist}). 

If $\Phi^G_{\beta,\nu}(\mu_1 \times \mu_2 | E(X,X')) \geq \gamma$, then on at least one of the pseudodistributions $\mu^\sym_1$ or $\mu^\sym_2$ Algorithm~\ref{alg:low-ent-restated} runs in time $\poly(|V(G)|)$ and returns an assignment of expected value at least 
\[
(\beta - \nu)^2\left(\gamma - O(\delta+\zeta) - \frac{1}{|V(G)|}\right) - 3\nu(\beta - \nu)
\]
for $I$.
\end{lemma}

%While Algorithm~\ref{alg:low-ent} is randomized, we can derandomize it and obtain a deterministic polynomial-time algorithm with the same guarantee on the approximation factor. To derandomize we can use standard techniques such as the method of conditional expectations~\cite{Vadhan12}. We will refer to such an algorithm as derandomized Condition \& Round.

 
\begin{proof}[Proof of Lemma~\ref{thm:round-j}]
Following the proof strategy of~\cite{BBKSS} we define the alternate shift potential $\Psi(\mu)$:
\begin{definition}[Alternate shift potential~\cite{BBKSS}]
The {\em alternate shift potential} of a degree-$D \ge 4$ pseudodistribution $\mu$ is given by
\[
\Psi(\mu) := \E_{u,v \sim \pi(G)} \left[\sum_{s \in \Sigma} \pPr_\mu[X_v - X_u = s]^2 \cdot \pE[\val_v(X) \mid X_v - X_u = s] \right],
\]
where $\pi(G)$ is the stationary measure on $G$ and $\val_v(X)$ denotes the value of vertex $v$.
%${\val_v(X) = \E_{w \sim v}[\pPr_\mu[X \text{ satisfies } (v,w)]]}$ for $w \sim v$ a  neighbor of $v$ sampled proportional to the weight on $(v,w)$.
\end{definition}

We will show that if $\Phi(\mu_1 \times \mu_2 | E)$ is large and $E$ does not introduce too many correlations, then either $\Psi(\mu_1)$ or $\Psi(\mu_2)$ must be large:
\begin{lemma}\label{lem:relating-ent-j}
Let $\mu_1 \times \mu_2$ be a degree $\widetilde{O}(D+1/\nu)$ pseudodistribution over $(X,X')$ and let $E(X,X')$ be a polynomial such that $\cA_I \vdash_D E(X,X') \in [0,1]$. Suppose we have that: 
\[\Pr_{u,v \sim V(G): u \neq v}[TV((Y_{u,v},Y'_{u,v})|E , (Y_{u,v},Y'_{u,v})) > \delta] \leq \zeta.\]

If the shift-partition potential of $\mu_1 \times \mu_2 |E(X,X')$ is large, then the alternate shift-potential of $\mu_1$ or $\mu_2$ must be large as well:
\[\Phi_{\beta,\nu}(\mu_1 \times \mu_2 | E(X,X')) \leq \frac{\Psi(\mu_1)}{2(\beta-\nu)^2} +\frac{\Psi(\mu_2)}{2(\beta-\nu)^2} + \frac{3\nu}{\beta - \nu} + 2\delta + 2\zeta + \frac{1}{|V(G)|}.\]
\end{lemma}
\begin{proof}
Deferred to Section~\ref{sec:lems}.
\end{proof}
Given a pseudodistribution $\mu$, we 
recall that $\mu^{sym}$ is the shift symmetrized pseudodistribution as in Section~\ref{sec:symm-prelims}.
We note that $\Psi(\mu) = \Psi(\mu^{\sym})$ because $\Psi$ as $\pPr_{\mu^\sym}[X_u - X_v = s] = \pPr_{\mu}[X_u - X_v = s]$ and $\pE_{\mu^\sym}[\val_v(X)|X_v - X_u = s] = \pE_{\mu}[\val_v(X)|X_v - X_u = s]$ for all $s$.  Therefore we can use the following lemma from~\cite{BBKSS} that shows that when the alternate shift potential of a shift-symmetric pseudodistribution is large, a single step of conditioning and rounding returns a solution of high objective value:
\begin{lemma}[Lemma 3.6~\cite{BBKSS}]
Let $I = (G, \Pi)$ be an affine instance of Unique Games over the alphabet $\Sigma$, let $\kappa>0$ and 
let $\mu$ be a degree-$4$ shift-symmetric pseudodistribution for $I$. If $\Psi(\mu) \geq \kappa$, then Algorithm~\ref{alg:low-ent-restated} returns a solution of expected value at least $\kappa$.
\end{lemma}

The statement now follows by applying the last lemma on either $\mu^\sym_1$ or $\mu^\sym_2$ (depending 
which one has a higher $\Psi$).
\end{proof}

 \subsubsection{Relating the potentials: Proof of Lemma~\ref{lem:relating-ent-j}}\label{sec:lems}
% In this section, we will prove Lemma~\ref{lem:relating-ent-j} and relate the potential $\Phi_{\beta,\nu}(\mu_1 \times \mu_2 | E)$ to $\Psi(\mu_1)$ and $\Psi(\mu_2)$. The proof is along the lines of Lemma~\ref{lem:relating-ent}, except that we now have to handle error terms that arise because we are working with a non-product pseudodistribution $(\mu_1 \times \mu_2 | E)$ over $(X,X')$.

% Before we dive into these lemmas, let us define a symmetrization operation on pseudodistributions. Intuitively it makes sense
% for a pseudodistribution on an affine unique games instance $I$ to be symmetric with respect to shifts, since if $X$ is a $(1-\eps)$-satisfying solution for $I$, then so is $X_{+s}$ for all $s \in \Sigma$. Pseudodistributions obtained by symmetrization will satisfy useful symmetry properties that are amenable to the analysis of Algorithm~\ref{alg:low-ent}.

% \begin{definition}[Symmetrization]\label{def:sym}
% Given a pseudodistribution $\mu$, we define the corresponding symmetrized pseudodistribution $\mu_{sym}$ as: For each $s \in \Sigma$, define $\mu_{+s}$ to be the pseudodistribution in which the labels receive the global affine shift $+s$, so that 
% \[
% \pE_{\mu_{+s}}[X_{u_1,a_1} \cdots X_{u_t,a_t}]
% := \pE_{\mu}[X_{u_1,a_1 - s}\cdots X_{u_t,a_t -s}]
% \]
% for all $\{(u_1,a_1), \ldots, (u_t,a_t)\} \in ([n] \times \Sigma)^{\le D}$. Now, define $\mu_{sym}$ to be the uniform mixture over $\mu_{+s}$ with $s \in \Sigma$. We say that a pseudodistribution is shift-symmetric if it is invariant under the symmetrization operation defined above, that is, $\mu = \mu_{sym}$.
% \end{definition}

% Firstly note that this operation can be efficiently performed on $\mu$. Furthermore it yields a valid pseudodistribution that has the same value as $\mu$.

% \begin{lemma}[Symmetrization]\label{lem:sym-basic}
% Let $\mu$ be a degree-$D$ pseudodistribution satisfying the unique games axioms $\A_I$ given by (\ref{eq:ip}) for an affine unique games instance $I$. Let $\mu_{sym}$ be a pseudoexpectation operator obtained by symmetrizing $\mu$. Then we have that,
% \begin{enumerate}
% \item $\mu_{sym}$ is a valid pseudoexpectation operator of degree-$D$ that satisfies the unique games axioms $\cA_I$. 
% \item The time taken to perform symmetrization on $\mu$ is subquadratic in the description of $\mu$.
% \item The objective value of $\mu$ and $\mu_{sym}$ are equal, i.e. $\val_\mu(I) = \val_{\mu_{sym}}(I)$.
% \end{enumerate}
% \end{lemma}
%The proof of this lemma is fairly straightforward, so we omit it. Since the value is invariant under symmetrization and performing the operation is efficient, all our algorithms symmetrize the pseudodistributions obtained by solving the degree $D$ SoS relaxation, and hence in our analysis we always work with shift-symmetric pseudodistributions henceforth.

% Symmetrized distributions satisfy some nice symmetry properties with respect to shifts, such as, every vertex has uniform marginals, and value of $\mu$ conditioned on $X_u = s$ for any shift $s$, is the same as the original value without conditioning. Additionally we have the following:

% \begin{lemma}[Shift-Symmetry properties]
% Let $\mu$ be a degree-$D$ shift-symmetric pseudodistribution satisfying the unique games axioms $\A_I$ given by (\ref{eq:ip}) for an affine unique games instance $I$.
% Then $\mu$ satisfies the following symmetry properties:
% \begin{enumerate}
%     \item For all vertices $u,v \in V(G)$ and shifts $s \in \Sigma$,
%     $\pPr[X_v = s \mid X_u = 0] = \pPr[X_v - X_u = s]$.
%     \item For all polynomials $f(X)$ with $\deg(f) \le D-2$, such that $f(X) = f(X+s)$ for every global shift $s \in \Sigma$,
%     \[\pE[f(X) \mid X_v-X_u=s] = \pE[f(X) \mid X_u = 0, X_v = s].\]
% \end{enumerate}
% \end{lemma}
% This lemma follows easily from the fact that $\mu$ is invariant under global shifts. See Appendix~\ref{sec:prelims} for a proof. 

% We first prove that when the potential $\Psi$ is large, conditioning and then independently rounding succeeds.

% \begin{proof}
% Suppose that $\Psi(\mu) \geq \delta$.
% Define the following,
% \[
% \Psi_{u}(\mu) := \E_{v \sim \pi} \left[\sum_{s \in \Sigma} \pPr[X_v - X_u = s]^2 \cdot \pE[\val_v(X) \mid X_v - X_u = s] \right],
% \]
% so that $\Psi(\mu) = \E_{u \sim \pi}[\Psi_u(\mu)]$. Suppose we sample a random vertex $u \sim \pi$ and condition the pseudodistribution on $X_u = 0$, then pick a random label $Y^u_v$ for every vertex $v \in V(G)$ according to its marginal $Y^u_v \sim \pE[X_v \mid X_u = 0]$.
% We have that in expectation, after conditioning on $u$ the rounded value is equal to:
% \[
% \E_{Y^u}[\val(Y^u)] = \E_{v\sim \pi}\E_{w\sim v}\left[\sum_s \pPr[X_v = s \mid X_u = 0]\pPr[X_w = \pi_{vw}(s) \mid X_u = 0] \right].\]

% We will now lower bound this quantity by $\Psi_u(\mu)$. 
% We have that
% \begin{align*}
% \Psi_{u}(\mu)
% &= \E_{v \sim \pi} \left[\sum_{s \in \Sigma} \pPr[X_v - X_u = s]^2 \cdot \pE[\val_v(X) \mid X_v - X_u = s] \right] \\
% &= \E_{v} \left[\sum_{s \in \Sigma} \pPr[X_v = s | X_u = 0]^2 \cdot \pE[\val_v(X) \mid X_u = 0, X_v = s] \right] 
% \intertext{where we have applied Lemma~\ref{lem:sym} along with the shift-symmetry of $\mu$ and of $\val_u(X)$, where the latter is a shift-symmetric function because $I$ is an affine unique games instance. Now, by definition of the local value,}
% &=\E_{v\sim \pi} \left[\sum_{s \in \Sigma} \pPr[X_v = s | X_u = 0]^2 \cdot \E_{w \sim v}\left[\pPr[X \text{ satisfies }(v,w) \mid X_u = 0, X_v = s]\right] \right] \\
% &= \E_{v \sim \pi} \left[\sum_{s \in \Sigma} \pPr[X_v = s | X_u = 0]^2 \cdot \E_{w \sim v}\left[\frac{\pPr[X_v = s, X_w = \pi_{vw}(s) \mid X_u = 0]}{\pPr[X_v = s \mid X_u = 0]}\right] \right] \\
% &= \E_{v\sim \pi} \left[\sum_{s \in \Sigma} \pPr[X_v = s | X_u = 0] \cdot \E_{w \sim v}\left[\pPr[X_v = s, X_w = \pi_{vw}(s) \mid X_u = 0]\right] \right] \\
% &\leq \E_{v \sim \pi} \E_{w \sim v} \left[\sum_{s \in \Sigma} \pPr[X_v = s | X_u = 0] \cdot \pPr[X_w = \pi_{vw}(s) \mid X_u = 0] \right] \\
% &= \E_{Y^u}[\val(Y^u)]
% \end{align*}
% Further note that the expected value of rounding of Algorithm~\ref{alg:low-ent} is $\E_{u \sim \pi}[\E_{Y^u}[\val(Y^u)]]$ which is greater than $\Psi(\mu)$ by the above inequality. Since $\Psi_{u}(\mu) \geq \delta$, we sample a solution with expected value at least $\delta$.
% \end{proof}
%Now, we will relate the two potentials.

\begin{proof}[Proof of Lemma~\ref{lem:relating-ent-j}]
We begin by recalling that in the definition of $\Phi_{\beta,\eta}$, we used an $\nu$-additive polynomial approximation $p(x)$ of degree $\widetilde{O}(1/\nu)$ to the indicator function $\Ind[x \ge \beta]$ on the interval $x \in [0,1]$, guaranteed by Theorem~\ref{thm:step-approx}. We will use $E$ to denote the polynomial $E(X,X')$, $\cD$ to denote the pseudodistribution $\mu_1 \times \mu_2 | E$ and $\pi$ to denote the uniform (in general stationary) measure over $V(G)$. Let $p_u$ denote the polynomial $p_{\beta,\nu}(\val_u(X))$ and $p'_u = p_{\beta,\nu}(\val_u(X'))$. Recall also the collection of local distributions $\cL(\cdot, p_{\beta,\nu})$ (Definition~\ref{defn:local-dist}). We will overload the notation $p_u,p'_u$ to mean the variables from Definition~\ref{defn:local-dist} and the polynomial $p_{\beta,\nu}(\val_u(X))$, and the use should be clear from context. Expanding the definition of $\Phi_{\beta,\nu}(\cD)$ we get:
\begin{align}
    \Phi_{\beta,\nu}(\cD)
    &= \pE_{\cD}\left[\sum_{s \in \Sigma} \left(\E_{u \sim \pi} \Ind[X_u - X_u' = s]p_u p'_u\right)^2\right ] \nonumber\\
    &= \pE_{\cD}\left[\sum_{s \in \Sigma} \E_{u,v \sim \pi} \Ind[X_u - X_u' = X_v - X_v' = s]p_up_v p'_u p'_v \right]\nonumber\\
    &= \E_{u,v \sim \pi}\pE_{\cD}\left[ \Ind[X_u - X_v = X'_u- X'_v]\cdot p_u p_vp'_up'_v\right], \nonumber\\
    &= \E_{u,v\sim \pi}[\pE_{\mu_1 \times \mu_2}\left[ \Ind[X_u - X_v = X'_u- X'_v]\cdot p_u p_vp'_up'_v\right] + {\sf err}(u,v)],\label{eq:relating-potential}
\end{align}
for ${\sf err}(u,v)$ defined accordingly for each pair $(u,v)$ (note that we have switched from the pseudo-distribution $\mathcal{D}$ to the pseudo-distribution $\mu_1\times \mu_2$ in the last transition).

\paragraph{Bounding the error term in~\eqref{eq:relating-potential}.} 
For $u = v$, we get that ${\sf err}(u,v) = \pE_{\cD}[p_u^2 (p'_u)^2] - \pE_{\mu_1 \times \mu_2}[p_u^2 (p'_u)^2] \leq 1$. Let us bound ${\sf err}(u,v)$ for $u \neq v$.
Let $R(X_u,X_v,X'_u,X'_v,p_{u},p_{v}, p'_u,p'_v)$ be the polynomial $\Ind[X_u - X_v = X'_u- X'_v]p_u p_v p'_u p'_v$ ($p_u$ denotes the polynomial $p_{\beta,\nu}(\val_u(X))$). Here we will overload the notation of $p_u$ to also denote the indicator variables from Definition~\ref{defn:local-dist} and let $R(a_1,\ldots,a_4,b_1,\ldots,b_4)$ for $a_1,\ldots,a_4 \in \Sigma$ and $b_1,\ldots,b_4 \in \{0,1\}$ also denote the evaluation of $R(X_u,\ldots,p'_v)$ on the values $a_1,\ldots,b_4$: $\Ind[a_1 - a_2 = a_3- a_4]b_1 b_2b_3 b_4$. Using Definition~\ref{defn:local-dist} we have that when $u \neq v$: $\pE_{\cD}[R] = \pPr_{\cL(\cD,p_{\beta,\nu})}[X_u - X_v = X'_u - X'_v, p_u = p_v = p'_u = p'_v = 1]$ and similarly for the distribution $\mu_1 \times \mu_2$ in place of $\cD$. Therefore we get that for $u \neq v$: 
\begin{align*}
&{\sf err}(u,v)\\
=&\pE_{\cD}[R(X_u,\ldots,p'_v)] - \pE_{\mu_1 \times \mu_2}[R(X_u,\ldots,p'_v)] \\
=&\pPr_{\cL(\cD,p_{\beta,\nu})}[X_u - X_v = X'_u - X'_v, p_u = \ldots =p'_v = 1] - \pPr_{\cL(\mu_1 \times \mu_2,p_{\beta,\nu})}[X_u - X_v = X'_u - X'_v, p_u = \ldots = p'_v = 1]\\
=&\sum_{\substack{a_1,\ldots,a_4 \in \Sigma}}  R(a_1,\ldots,a_4, 1,\ldots,1)(\pPr_{\cL(\cD,p_{\beta,\nu})}[X_u = a_1,\ldots,p_u =1, \ldots] - \pPr_{\cL(\mu_1 \times \mu_2,p_{\beta,\nu})}[X_u = a_1,\ldots,p_u =1 \ldots])\\
&\leq \sum_{\substack{a_1,\ldots,a_4 \in \Sigma}} \left|\pPr_{\cL(\cD,p_{\beta,\nu})}[X_u = a_1,\ldots,p_u = 1,\ldots] - \pPr_{\cL(\mu_1 \times \mu_2,p_{\beta,\nu})}[X_u = a_1,\ldots,p_u = 1,\ldots] \right| \\
&= 2 TV((Y_{u,v},Y'_{u,v})|E , (Y_{u,v},Y'_{u,v})),
\end{align*}
where $Y_{u,v} = (X_u,X_v,p_u,p_v)$, $Y'_{u,v} = (X'_u,X'_v,p'_u,p'_v)$ as in the lemma, the distribution $(Y_{u,v},Y'_{u,v})|E$ refers to the joint distribution on these variables defined by the collection of local distributions $\cL(\mu_1 \times \mu_2|E, p_{\beta,\nu})$ and similarly $(Y_{u,v},Y'_{u,v})$ refers to the distribution defined by $\cL(\mu_1 \times \mu_2,p_{\beta,\nu})$ (Definition~\ref{defn:local-dist}). 

Combining, we get that: 
\begin{align}\label{eq:err-uv}
\E_{u,v \sim \pi}[{\sf err}(u,v)] &\leq \E_{u,v \sim \pi: u \neq v}[{\sf err}(u,v)] + \Pr_{u, v \sim \pi}[u = v]\E_{u}[{\sf err}(u,u)] \notag\\
&\leq \E_{u,v \sim \pi: u \neq v}[2 TV((Y_{u,v},Y'_{u,v})|E , (Y_{u,v},Y'_{u,v}))] + \frac{1}{|V(G)|} \notag\\
&\leq 2\delta + 2\zeta + \frac{1}{|V(G)|},
\end{align}
by the assumption in the lemma statement. 
\paragraph{Bounding the main term in~\eqref{eq:relating-potential}.} We now upper bound the first term in~\eqref{eq:relating-potential}. 
\begin{align*}
&\E_{u,v\sim \pi}[\pE_{\mu_1 \times \mu_{2}}\left[ \Ind[X_u - X_v = X'_u- X'_v]\cdot p_u p_vp'_up'_v\right] \\
&=\E_{u,v\sim \pi}\left[\sum_{s \in \Sigma} \pE_{\mu_1}[\Ind[X_u - X_v = s]\cdot p_u p_v]\pE_{\mu_2}[\Ind[X'_u - X'_v = s]\cdot p'_u p'_v]\right].
\end{align*}

We next use that $p(x),\val_u(X),\val_u(X') \in [0,1]$ and Fact~\ref{fact:bdd-markov} asserting that 
$p(x) \le \frac{x}{\beta-\nu} + \nu$ for all $x \in [0,1]$, and furthermore this is \sos-certifiable.
Thus, pulling out a factor of $p_v,p'_v$ and applying
Fact~\ref{fact:bdd-markov} to bound $p_u,p'_u$ we get 
that the first term on the right hand side of~\eqref{eq:relating-potential} is at most
\begin{align*}
&\E_{u,v\sim \pi}\left[\sum_{s \in \Sigma} \pE_{\mu_1}\left[\Ind[X_u - X_v = s]\cdot \frac{\val_u(X)}{\beta - \nu}\right]\pE_{\mu_{2}}\left[\Ind[X'_u - X'_v = s]\cdot \frac{\val_u(X')}{\beta - \nu}\right]\right] + \frac{3\nu}{\beta-\nu} \\
&\leq \frac{1}{2}\E_{u,v\sim \pi}\left[\sum_{s \in \Sigma} \pE_{\mu_1}\left[\Ind[X_u - X_v = s]\cdot \frac{\val_u(X)}{\beta - \nu}\right]^2\right] + \frac{1}{2}\E_{u,v\sim \pi}\left[\sum_{s \in \Sigma}\pE_{\mu_{2}}\left[\Ind[X'_u - X'_v = s]\cdot \frac{\val_u(X')}{\beta - \nu} \right]^2\right] \\
&+ \frac{3\nu}{\beta-\nu},
\end{align*}
where we have used that the AM-GM inequality. The first term may be bounded as
\begin{align*}
&\E_{u,v\sim \pi}\left[\sum_{s \in \Sigma} \pE_{\mu_1}\left[\Ind[X_u - X_v = s]\cdot \frac{\val_u(X)}{\beta - \nu}\right]^2\right] \nonumber\\
\leq&\E_{u,v\sim \pi}\left[\sum_{s \in \Sigma} \pE_{\mu_1}\left[\Ind[X_u - X_v = s]\cdot \frac{1}{\beta-\nu}\right] \pE_{\mu_1}\left[\Ind[X_u - X_v = s]\frac{\val_u(X)}{\beta - \nu}\right]\right] \nonumber \\
=&\frac{1}{(\beta-\nu)^2}\E_{u,v\sim \pi}\left[ \sum_{s \in \Sigma} \pE_{\mu_1}[\Ind[X_u - X_v = s]]^2 \cdot \pE_{\mu_1}[\val_u(X) \mid X_u - X_v = s] \right]\nonumber\\
=&\frac{\Psi(\mu_1)}{(\beta-\nu)^2},
\end{align*}
where we have used that $\val_u(X) \leq 1$ and applied the definition of conditional pseudoexpectation, 
and similarly the second term is upper bounded by $\frac{\Psi(\mu_{2})}{(\beta-\nu)^2}$. Plugging this and~\eqref{eq:err-uv} into~\eqref{eq:relating-potential} finishes the proof.
\end{proof}

% \subsection{Iteration}

% \begin{lemma}\label{lem:subroutine}
% Let $\eps_0 \in (0,1)$ be a universal constant and $\delta: [0,1] \rightarrow [0,1]$ be a function. Let $\eps < \eps_0/2$ be any constant, and let $\delta_{\min} = \min_{\eta \in [\eps,2\eps]}(\delta(\eta))$. 
% Let $G$ be a regular graph and $I$ be any unique games instance on $G$ with alphabet size $|\Sigma| = k \geq \Omega(\frac{1}{\delta_{\min}})$ and value $1 - \eps$. 

% Suppose we have a subroutine $\cA$ which given as input $\mu$, a shift-symmetric degree-$D$ pseudodistribution satisfying $\cA_I$ with $\val_\mu(I) = 1 - \eta \ge 1-\eps_0$, returns a vertex-induced subgraph $H$ such that, 1) $\text{Sub-Round-val}_{\mu}(H) \geq \delta = \delta(\eta)$ and 2) the edge-expansion of $H$ is $O(\eta)$. 

% Then if $\cA$ runs in time $T(\cA)$, there is a $|V(G)|(T(\cA) + |V(G)|^{O(D)})$-time algorithm which finds a solution for $I$ that satisfies an $\Omega(\delta_{\min}^2\eps)$-fraction of the  edges of $G$. 
% \end{lemma}

% \begin{proof}
% We will use the algorithm $\cA$ as a subroutine. To get a full assignment, our algorithm below is a generalized version of the Algorithm~\ref{alg:j}, where we've replaced the steps 2 to 4 in Algorithm~\ref{alg:j} with an arbitrary subroutine $\cA$ that finds a subgraph with high Condition\&Round value with respect to $I$. We include it here for completeness.

% \begin{algorithm-thm}[Partial to Full Assignment]\label{alg:partial} ~\\ 
% \begin{compactenum}
% \item Solve the degree-$D$ SoS SDP relaxation for the integer program $\cA_I$ and make the pseudodistribution shift-symmetric to get a pseudodistribution $\mu_0$. Set $j = 1$.

% \item While $\eta_{j-1} := 1 - \val_{\mu_{j-1}}(I) \leq 2\eps$:
% \begin{compactenum}
% \item Run subroutine $\cA$ on $\mu_{j-1}$ to find a subgraph $H_{j}$ with $\text{CR-val}_{\mu_{j-1}}(H_j) \geq \delta(\eta_{j-1})$.
% \item Let $S_j$ be a subgraph of $H_j$ induced by the set of vertices that have not been previously assigned by any partial assignment $f_{k}, k < j$. Perform derandomized Condition\&Round on $V(S_j)$ to get a partial assignment $f_j$.
% \item Rerandomize the pseudodistribution $\mu_{j-1}$ on $S_j$ to get $\mu_j$: Make the marginal distribution over the assigned vertices uniform and independent of other vertices, that is, for all degree $\leq D$ monomials define
% $\pE_{\mu_{j}}$ as follows, 
% \[\pE_{\mu_{j}}[X_{h_1,a_1} \cdots X_{h_t,a_t}X_{u_1,b_1} \cdots X_{u_m,b_m}]
% := \frac{1}{|\Sigma|^t}\pE_{\mu_{j-1}}[X_{u_1,b_1}\cdots X_{u_m,b_m}],
% \]
% where $\{(h_1,a_1), \ldots, (h_t,a_t)\} \in (V(S_j) \times \Sigma)^{t}$ and $\{(u_1,b_1), \ldots, (u_m,b_m)\} \in (({[n] \choose \ell} \setminus V(S_j)) \times \Sigma)^{m}$.
% \item Increment $j$.
% \end{compactenum}
% \item Output any assignment to $V(G)$ that agrees with all partial assignments $f_j$ considered above.
% \end{compactenum}

% \end{algorithm-thm}


% Let us first check that the algorithm is well-defined. The initial pseudodistribution $\mu_0$ by definition satisfies axioms $\cA_I$ and is shift-symmetric. It has value $= 1 - \eps > 1 - \eps_0$. In subsequent iterations of the while loop all these properties are satisfied by the modified pseudodistributions: 1) the rerandomizing operation on pseudodistributions produces a valid pseudodistribution operator that satisfies the axioms $\cA_I$ and is also shift-symmetric, 2) At iteration $j$ of the while-loop, since the while condition is met, we know that $\mu_{j-1}$ has value $\geq 1 - 2\eps \geq 1 - \eps_0$ and furthermore we can show that since the value only decreases at each step, it is always $\leq 1 - \eps$, so that $\eta_j \in [\eps,2\eps]$. So inside the while-loop, $\cA$ will always find a non-empty subgraph $H_j$ with high Condition\&Round value. Next, we find an assignment $f_j$ to the set of vertices $V(S_j)$ that by definition don't intersect previously assigned vertices. Since $f_j$ doesn't reassign any vertices, in the final step of the algorithm it is possible to output an assignment that is consistent with all previously considered partial assignments. We will now show that our final partial assignment satisfies a large fraction of the edges, where we say that an edge $(u,v)$ is satisfied by a partial assignment $f_j$, if both vertices $u,v$ have been assigned labels under $f_j$ and the labels satisfy the edge. We claim the following two facts:

% \begin{claim}
% The drop in value in every iteration satisfies that:
% \[\val_{\mu_{j-1}}(I) - \val_{\mu_j}(I) \leq \frac{2|V(H_j)|}{|V(G)|},\]

% where $\val_\mu(I)$ denotes the SDP value of $I$ with respect to the pseudodistribution $\mu$. 
% \end{claim}

% \begin{proof}
% For any edge $(h,v)$ where $h \in H_j$, we have that $\val_{\mu_j}((h,v)) = \frac{1}{k}$, whereas $\val_{\mu_{j-1}}((h,v)) \leq 1$. For any edge whose both endpoints lie outside $H_j$, the value remains unchanged under rerandomizing. Noting that the fraction of edges incident on vertices in $H_j$ is at most $\frac{2|V(H_j)|}{|V(G)|}$ the conclusion follows.
% \end{proof}


% \begin{claim}
% The value of the partial assignment found at iteration $j$ satisfies:
% \[\val(f_j) \geq \delta^2(\eta_{j-1})(1 - O(\eta_{j-1}))\frac{|V(H_j)|}{2|V(G)|},\]
% where $\val(f_j)$ denotes the fraction of edges (in $E(G)$) satisfied by the partial assignment $f_j$.
% \end{claim}

% \begin{proof}
% We will first prove that $\text{CR-val}_{\mu_{j-1}}(S_j) \geq \delta(\eta_{j-1}) := \delta$, where $S_j$ is the subgraph induced by the unassigned (by previous partial assignments $f_k$, $k < j$) vertices of $H_j$. For notational simplicity we will drop the subscript $j$ from $H_j,S_j$ and $\mu_{j-1}$. We know by the guarantees of the subroutine $\cA$ that $H$ is such that, $\text{CR-val}_{\mu}(H) = \E_{u \sim V(H)}[\text{ind-val}_{\mu | X_u = 0}(H)] \geq \delta$ (see Definition~\ref{def:cr-val} for $\text{CR-val}$ and $\text{ind-val}$).  First note that the marginals of every vertex are uniform, due to the shift-symmetry of $\mu$. Moreover we have that conditioning on previously assigned vertices, i.e. any vertex $u \in H \setminus S$, maintains this property, since the distribution of $u$ is completely uncorrelated with the other vertices. So we get that, $\text{ind-val}_{\mu | X_u = 0}(H) = \frac{1}{|\Sigma|} < \delta$ for all $u \in H \setminus S$. This implies that, 
% \begin{equation}\label{eq:cr1}
% \E_{u \sim S}[\text{ind-val}_{\mu | X_u = 0}(H)] \geq \text{CR-val}_{\mu}(H) \geq \delta.    
% \end{equation}
% Again we have that, $\text{ind-val}_{\mu | X_u = 0}(e) = \frac{1}{|\Sigma|}$, for any edge $e$ which has at least one endpoint in $H \setminus S$, so we get that, $\E_{u \sim S}[\text{ind-val}_{\mu | X_u = 0}(S)] \geq \delta$. Now we can perform derandomized Condition\&Round on $S$ to get an assignment $f_j$ that satisfies at least a $\delta$-fraction of the edges of $S$.

% Next we will show, by an averaging argument, that the edges of $S$ constitute a large fraction of the edges incident on the vertices of $H$, which would imply that $f_j$ satisfies a large fraction of these edges. Let $u_0 \in S$ be a vertex for which $\text{ind-val}_{\mu | X_{u_0} = 0}(H) \geq \delta$ (we know such a vertex exists by equation~(\ref{eq:cr1})) and let $\mu'$ be the pseudodistribution $(\mu | X_{u_0} = 0)$. First note that the set of edges $E(H) \setminus E(S)$, have independent rounding value $1/|\Sigma|$ under $\mu'$, since at least one endpoint of such edges has been assigned previously. Since $\text{ind-val}_{\mu'}(H) \geq \delta$, a simple averaging argument gives us that the set $E(H) \setminus E(S)$ can be at most a $\frac{1 - \delta}{1 - (1/|\Sigma|)}$-fraction of $E(H)$. So the set $E(S)$ is at least a $\frac{\delta - (1/|\Sigma|)}{1 - (1/|\Sigma|)} \geq \delta/2$-fraction of $E(H)$. Since the expansion of $H$ is at most $O(\eta_{j-1})$, we have that $E(H)$ is a $(1 - O(\eta_{j-1}))$-fraction of the total edges incident on $H$, which in turn is at least a $\frac{|V(H)|}{|V(G)|}$-fraction of $E(G)$. Combining these facts we get that $f_j$ satisfies a $\delta \cdot \frac{\delta}{2} \cdot (1 - O(\eta_{j-1}))\cdot \frac{|V(H)|}{|V(G)|}$-fraction of the edges of $G$. 
% \end{proof}

% Once we have these facts, the conclusion is immediate. Firstly there cannot be more than $V(G)$ iterations of the while-loop, since  at each iteration we set the value of at least one new vertex to $1/|\Sigma|$. The rerandomization operation in the while loop as well as the symmetrization operation (Lemma~\ref{lem:sym}) can be done in time polynomial in the description of $\mu_0$. So each iteration takes time $T(\cA) + |V(G)|^{O(D)}$, hence the algorithm runs in time $|V(G)|(T(\cA) + |V(G)|^{O(D)})$. 

% Moreover, combining the claims above, we get that the partial assignment at any iteration is proportional to the drop in value of the pseudodistribution. That is, 

% \[\val(f_j) \geq \delta^2(\eta_{j-1})(1 - O(\eta_{j-1}))\frac{|V(H_j)|}{2|V(G)|} \geq \delta^2(\eta_{j-1})(1 - O(\eta_{j-1}))\left(\frac{\val_{\mu_{j-1}}(I) - \val_{\mu_j}(I)}{4}\right).\]

% At the last iteration, we know that the pseudodistribution value has dropped by at least $\eps$ (compared to $\val_{\mu_0}(I)$), hence summing the above over all iterations $j$, we get that the value of the partial assignment returned by the algorithm is at least ${\frac{1}{4}(\min_j(\delta(\eta_j)))^2(1 - O(\eps))\eps}$ as required.
% \end{proof}




